\begin{Abstrak}
    Produksi literatur ilmiah di lingkungan perguruan tinggi menunjukkan peningkatan yang signifikan setiap tahunnya. Peningkatan jumlah publikasi menyebabkan peneliti mengalami kesulitan dalam menemukan literatur yang relevan secara efisien. Sistem pencarian akademik konvensional yang mengandalkan pencocokan kata kunci kurang mampu memahami konteks semantik dan relasi antar-entitas akademik. Penelitian ini bertujuan merancang \textit{pipeline} konstruksi \textit{Knowledge Graph} dari korpus literatur akademik, merancang mekanisme \textit{Hybrid Vector-Graph Retrieval} dalam arsitektur \textit{Knowledge Graph-Based Retrieval Augmented Generation}, dan mengevaluasi efektivitas integrasi \textit{Knowledge Graph} dalam meningkatkan kualitas jawaban model bahasa besar. Data publikasi diperoleh dari profil dosen dan publikasi rumpun Informatika dan Komputer UNESA. Konstruksi graf menghasilkan 45.685 simpul dan 122.675 relasi dengan 16 jenis simpul yang mencakup entitas bibliografis dan konseptual. Eksperimen dilakukan dengan membandingkan tiga metode \textit{retrieval}, yaitu \textit{Naive}, \textit{Subgraph}, dan \textit{Hybrid} pada 90 pertanyaan evaluasi yang terbagi ke dalam kategori faktual, relasional, dan \textit{multi-hop}. Evaluasi \textit{retrieval} menggunakan \textit{Mean Reciprocal Rank} dan \textit{Hit@K}, sedangkan evaluasi generasi jawaban menggunakan metode \textit{Pairwise LLM as Judge}. Metode \textit{Hybrid} menghasilkan performa terbaik dengan \textit{MRR} sebesar 0,50 dan \textit{Win Rate Overall} sebesar 53,9\% terhadap metode \textit{Naive}. Metode \textit{Subgraph} memperoleh \textit{MRR} sebesar 0,43 dengan \textit{Win Rate Overall} sebesar 45,3\% terhadap \textit{Naive}. Evaluasi pada dimensi \textit{Traceability} dan \textit{Faithfulness} menunjukkan integrasi \textit{Knowledge Graph} membantu model bahasa besar menghasilkan jawaban yang lebih mudah dilacak sumbernya dan mengurangi klaim yang tidak terverifikasi.

    \katakunci{\textit{Knowledge Graph}, \textit{Retrieval Augmented Generation}, \textit{Hybrid Retrieval}, Pencarian Literatur Akademik, Model Bahasa Besar}
\end{Abstrak}
